%Dokumentklasse

%draft als optionohne bilder für bessere performance
%\documentclass[a4paper,12pt,draft]{scrreprt}

%normal mit Bildern
\documentclass[a4paper,12pt]{scrreprt}

\usepackage[left= 3cm,right = 3cm, bottom = 3cm,top = 3cm]{geometry}
%\usepackage[onehalfspacing]{setspace}

% ============= Packages =============
% Dokumentinformationen
\usepackage[
pdftitle={Praktikum - Umwelttechnik},
pdfsubject={},
pdfauthor={Roman-Luca Zank},
pdfkeywords={},	
%Links nicht einrahmen
hidelinks
]{hyperref}

%nur Text zum prüfen des Umfangs

% Standard Packages
%\usepackage[bottom]{footmisc}
\usepackage[utf8]{inputenc}
\usepackage[ngerman]{babel}

\usepackage[T1]{fontenc}
%\usepackage{helvet}

%\renewcommand{\familydefault}{\sfdefault}

\usepackage{graphicx}
\graphicspath{{img/}}
\usepackage{mhchem}
\usepackage{fancyhdr}
\usepackage{lmodern}
\usepackage{color}
\usepackage{placeins}
\usepackage{booktabs}
\usepackage{caption}
\usepackage[list=true]{subcaption}
\usepackage{longtable}
\usepackage{tikz}
\usepackage{pgfplots}
\usepackage{lastpage}
\usetikzlibrary{patterns}

%Einheitenpackage
\usepackage{siunitx}  
\sisetup{	locale = DE, 
	per-mode=fraction,
	inter-unit-product=\ensuremath{\cdot},
	detect-weight = true,
	quotient-mode=fraction
}
%neue Einheiten definieren
\DeclareSIUnit\xyz{xyz}	
\DeclareSIUnit\rpm{rpm}	

%Automatisch cdot statt *
\DeclareMathSymbol{*}{\mathbin}{symbols}{"01}



%\setcounter{lofdepth}{2}      %für subfigures in list of figures

%Tabelle
\usepackage{tabularx}
\usepackage{tabulary}

%%Verzeichnispackages
%\usepackage{natbib}

%\usepackage{hyperref}
%\newcites{

% zusätzliche Schriftzeichen der American Mathematical Society
\usepackage{amsfonts}
\usepackage{amsmath}

%Abkürzungsverzeichnis
\usepackage{acronym}

%kein Abstand bei neuem Kapitel vom Seitenanfang
%\vspace*{2.3\baselineskip} = ORIGINAL
\renewcommand*{\chapterheadstartvskip}{\vspace*{.0\baselineskip}}

%nicht einrücken nach Absatz
\setlength{\parindent}{0pt}


\urlstyle{same}

% ============= Kopf- und Fußzeile =============
\pagestyle{fancy}
%
\lhead{}
\chead{}
\rhead{}%\slshape }%\leftmark}
%%
\lfoot{}
\cfoot{}
\rfoot[{\thepage\ of \pageref*{LastPage}}]{Seite \thepage\ von \pageref*{LastPage}}
%%
\renewcommand{\headrulewidth}{0pt}
\renewcommand{\footrulewidth}{0pt}
\renewcommand{\chapterpagestyle}{fancy}

%Fußnotelinie
%\let\footnoterule

%Fußnote mit Klammer
\renewcommand*{\thefootnote}{(\arabic{footnote})}

%Abb. statt Abbildung
\addto\captionsngerman{%
	\renewcommand{\figurename}{Abb.}%
	\renewcommand{\tablename}{Tab.}%
}

% ============= Package Einstellungen & Sonstiges ============= 
%Besondere Trennungen
%\hyphenation{De-zi-mal-tren-nung}
\usepackage[none]{hyphenat}
\hyphenpenalty=5000
\tolerance=5000
\providecommand\phantomsection{}

\usepackage{mathtools}


% ============= Dokumentbeginn =============

\begin{document}
%Seiten ohne Kopf- und Fußzeile sowie Seitenzahl
\pagestyle{empty}

%\begin{center}
\begin{tabular}{p{\textwidth}}


\begin{center}
\includegraphics[scale=0.75]{logos.jpg}\\
\end{center}


\\

\begin{center}
\LARGE{\textsc{
Protokoll \\
Umwelttechnik\\
}}
\end{center}

\\

%\begin{center}
%\large{Fakultät für Muster und Beispiele \\
%der Hochschule Musterhausen \\}
%\end{center}
%
%\\

\begin{center}
\textbf{\Large{V3 - Abwasserbeschaffenheit II}}
\end{center}

\begin{center}
	\large{Gruppe 1.2 (BCUT3)}
\end{center}


\\
%\begin{center}
%zur Erlangung des akademischen Grades\\
%Bachelor of Engineering
%\end{center}


%\begin{center}
%vorgelegt von
%\end{center}

\begin{center}
\Large{\textbf{Teilnehmer:}} \\ 
\end{center}
\begin{center}
\large{Christoph Hecht \\
	Willy Messerschmidt \\
	Roman-Luca Zank} \\
\end{center}


\\

\begin{center}
\begin{tabular}{lll}
\large{\textbf{Protokollführer:}} & & \large{Roman-Luca Zank} \\
& & \href{mailto:roman-luca.zank@stud.hs-merseburg.de}{{\footnotesize roman-luca.zank@stud.hs-merseburg.de}}\\
&&\\
\large{\textbf{Datum der Versuchsdurchführung:}}&& \large{17.12.2019}\\
&&\\
\large{\textbf{Abgabedatum:}}&& \large{XX.12.2019}
\end{tabular}
\end{center}

\\ \\ \\ \\ \\ \\ 
\large{Merseburg den \today}

\end{tabular}
\end{center}


%\include{14_danksagungen}

%\include{15_zusammenfassung}

% Beendet eine Seite und erzwingt auf den nachfolgenden Seiten die Ausgabe aller Gleitobjekte (z.B. Abbildungen), die bislang definiert, aber noch nicht ausgegeben wurden. Dieser Befehl fügt, falls nötig, eine leere Seite ein, sodaß die nächste Seite nach den Gleitobjekten eine ungerade Seitennummer hat. 
\cleardoubleoddpage

% Pagestyle für Titelblatt leer
\pagestyle{empty}

%Seite zählen ab
\setcounter{page}{0}

%Titelblatt
\begin{center}
\begin{tabular}{p{\textwidth}}


\begin{center}
\includegraphics[scale=0.75]{logos.jpg}\\
\end{center}


\\

\begin{center}
\LARGE{\textsc{
Protokoll \\
Umwelttechnik\\
}}
\end{center}

\\

%\begin{center}
%\large{Fakultät für Muster und Beispiele \\
%der Hochschule Musterhausen \\}
%\end{center}
%
%\\

\begin{center}
\textbf{\Large{V3 - Abwasserbeschaffenheit II}}
\end{center}

\begin{center}
	\large{Gruppe 1.2 (BCUT3)}
\end{center}


\\
%\begin{center}
%zur Erlangung des akademischen Grades\\
%Bachelor of Engineering
%\end{center}


%\begin{center}
%vorgelegt von
%\end{center}

\begin{center}
\Large{\textbf{Teilnehmer:}} \\ 
\end{center}
\begin{center}
\large{Christoph Hecht \\
	Willy Messerschmidt \\
	Roman-Luca Zank} \\
\end{center}


\\

\begin{center}
\begin{tabular}{lll}
\large{\textbf{Protokollführer:}} & & \large{Roman-Luca Zank} \\
& & \href{mailto:roman-luca.zank@stud.hs-merseburg.de}{{\footnotesize roman-luca.zank@stud.hs-merseburg.de}}\\
&&\\
\large{\textbf{Datum der Versuchsdurchführung:}}&& \large{17.12.2019}\\
&&\\
\large{\textbf{Abgabedatum:}}&& \large{XX.12.2019}
\end{tabular}
\end{center}

\\ \\ \\ \\ \\ \\ 
\large{Merseburg den \today}

\end{tabular}
\end{center}
 %Prokolle
%\begin{center}
\begin{tabular}{p{\textwidth}}


\begin{center}
\includegraphics[scale=0.75]{img/logos.jpg}\\
\end{center}


\\

\begin{center}
\LARGE{\textsc{
Recherche \\
Rückgewinnung von Ammoniak aus Industrieabwässern\\
}}
\end{center}

%\begin{center}
%\large{Fakultät für Muster und Beispiele \\
%der Hochschule Musterhausen \\}
%\end{center}
%
%\\
 \\
 
\begin{center}
\textbf{\Large{Seminararbeit in Medienrecherche}}
\end{center}

\begin{center}
	\large{im WiSe 2019}
\end{center}
 \\
%\begin{center}
%zur Erlangung des akademischen Grades\\
%Bachelor of Engineering
%\end{center}


\begin{center}
\large{vorgelegt von}
\end{center}
\\


\begin{center}
\Large{\textbf{Roman-Luca Zank}} \\
\end{center}

\begin{center}
3. Semester \\
Chemie- und Umwelttechnik \\
\end{center}


\begin{center}
\begin{tabular}{lll}
	\textbf{E-Mail:} & & romanzank@mail.de\\
	\textbf{Matrikelnummer:} & &25240\\
	\textbf{Adresse:} & &Platz der Bausoldaten 2, Zimmer 224\\
	\textbf{Ort:} & &06217 Merseburg\\
	&& \\
	\textbf{Prüfer:} & & Dr. Frank  Baumann\\
\end{tabular}
\end{center}

\\ \\ \\ \\ \\
\large{Merseburg, \today}

\end{tabular}
\end{center}
 %Seminar-/Abschlussarbeit

% Pagestyle für Rest des Dokuments
\pagestyle{fancy}

%Inhaltsverzeichnis
\tableofcontents
\thispagestyle{empty}

%Inhalt

%%Verzeichnis aller Bilder
\label{sec:bilder}
\listoffigures
\addcontentsline{toc}{chapter}{Abbildungsverzeichnis}
\thispagestyle{empty}

%Verzeichnis aller Tabellen
\label{sec:tabellen}
\listoftables
\addcontentsline{toc}{chapter}{Tabellenverzeichnis}
\thispagestyle{empty}

%Abkürzungsverzeichnis
\setlength{\columnsep}{20pt}
%\twocolumn
%\addchap{Nomenklatur}
%\label{sec:abkurzung}
%\begin{acronym}
%\acro{kf}[$\text{k}_\text{f}$]{Durchlässigkeitsbeiwert}
%\acro{t}{Durchlaufzeit}
%\acro{tm}[$\text{t}_\text{m}$]{Mittlere Durchlaufzeit}
%\acro{V}{Volumen}
%\acro{h}{Höhe der Wassersäule}
%\acro{Q}{Volumenstrom}
%\acro{l}{Durchströmte Länge}
%\acro{A}{Grundfläche}
%\acro{d}{Durchmesser}
%
%\end{acronym}
%\subsubsection{Aufrufen einer Abkürzung}
%\acs{rT}
%\begin{verbatim}
%\acs{Abkürzung}
%\end{verbatim}

\chapter{Aufgabenstellung}
\label{sec:aufgabenstellung}
%In der Aufgabenstellung wird (in eigenen Worten und ganzen Sätzen) formuliert, was das Ziel des 
%Versuches ist.  
%[Beachten Sie die eigentliche Aufgabenstellung in den Versuchsanleitungen sowie die Hinweise zur Auswertung!] 

Im Versuch 3 "`Abwasserbeschaffenheit II"' werden drei Abwasserproben unbekannter Herkunft ... \\
Ziel der Auswertung, der gesammelten Messdaten, ist eine Einschätzung der Herkunft und Belastung der Abwasserproben, sowie ein Vergleich der jeweiligen Beschaffenheit mit häuslichem Abwasser und den Mindestanforderungen für das Einleiten kommunaler Abwässer in einen Vorfluter der GK 5. \\
Im Anschluss sind Empfehlungen zur Abwasserbehandlung zu geben.

\chapter{Geräte und Chemikalien}
\label{sec:geraete}

\textbf{Geräte:}
\begin{itemize}
	\item Magnetrührer mit Rührfisch
	\item Bechergläser
	\item Erlenmeyerkolben
	\item Filterpapier
	\item \textsc{Imhoff}-Trichter (\SI{1}{\liter})
	\item Stoppuhr
	\item Trocknungswaage
	\item Muffelofen
	\item CSB-Reaktor mit Ampullen
	\item Kolorimetrisches Messgerät für CSB
	\item BSB-Win2000 von \textsc{Nanotec} (Differenzdruckmethode)
		\begin{itemize}
			\item Gerätesockel mit Messeinrichtung und Rührflächen
			\item volumengeeichte BSB$_5$–Flaschen
			\item computergestützte Messwerterfassung 
		\end{itemize}
	
\end{itemize}

\vspace*{5mm}

\textbf{Proben/Chemikalien:}
\begin{itemize}
	\item destilliertes Wasser
	\item Abwasserproben 1, 2 \& 3
	\item Kaliumdichromat-Lösung \textcolor{red}{haben wir das echt benutzt oder nicht das Schwefelzeug ?}
	\item \SI{1}{\milli \liter} vorbereitete Nährstoff-Ersatzlösung für \SI{300}{\milli \liter} Probe
	\item 2 x Dosiereinheiten Nitrifikationshemmer für \SI{300}{\milli \liter} Probe
	\item Schlifffett
	\item 2 x \ce{NaOH}-Plätzchen
\end{itemize}
%Start





\chapter{Durchführung}
\label{sec:durchfuerung}
Im ersten Versuchsteil werden die drei Proben mittels Abwasser-Feststoffuntersuchungen, sprich über das Sedimentationsverhalten, das Absetzvolumen (AV), die abfiltrierbaren Stoffe in Form der Trockensubstanz (TS) und der organischen Trockensubstanz (oTS), analysiert. In allen Versuchsteilabschnitten werden die Proben zuvor durch Schütteln homogenisiert.\\
Die Bestimmung des Sedimentationsverhaltens erfolgt durch einfüllen der drei Proben in je einen Imhoff-Trichter. Es wird Periodisch die Lage der Trennschicht abgelesen, welche sich zwischen der klaren und der Schwebstoffphase ausbildet.
Die Werte wurden 10 mal minütlich und danach alle 5 Minuten aufgenommen. Der letzte Wert wird nach insgesamt 30 Minuten abgelesen und ist gleichzeitig das Absetzvolumen. \\
Die Trockensubstanz wird durch Abfiltrieren einer bestimmten Probenmenge bestimmt. Vorher muss das verwendete Filterpaper getrocknet und eingewogen werden. Der erhaltene Filterkuchen wird auf dem Filterpapier bei \SI{105}{\degreeCelsius} getrocknet und zuletzt gewogen.\\
Der organische Anteil der Trockensubstanz erfolgt analog der DIN 38409. Der zuvor erhaltene Filterkuchen wird 2 Stunden lang bei \SI{550}{\degreeCelsius} verascht.\\

Im zweiten Versuchsteil werden die Proben über die Abwasser-Summenparameter in Form des chemischen Sauerstoffbedarfs (CSB) und des biochemischen Sauerstoffbedarfs über 5 Tage (BSB$_5$) untersucht.\\

Der chemische Sauerstoffbedarf wird mittels CSB-Testampulle bestimmt. Die Ampullen wurden mit den zu untersuchenden Proben befüllt, geschüttelt und in den vorgewärmten CSB-Reaktor zur Inkubation eingelegt und darin für 2 Stunden belassen. Auf den Deckeln der Proben ist der Inhalt zu beschriften. Es muss eine sinnvolle Verdünnung der Probe gewählt werden. In diesem Falle wurde für die Probe 1 ein hoher CSB erwartet. Darum wurde die Probe 1 ein mal im Verhältnis 1:1 mit Wasser gemischt, und ein mal pur analysiert. Die eigentliche Analyse wurde mit Hilfe des Hach-Messgerätes entsprechend der Betriebsanweisung durchgeführt.\\

Die Vorbereitung der Proben für die Ermittelung des Biochemischen Sauerstoffbedarfes begann mit der Befüllung der Analyseflaschen. In diese wurden je ein Magnetrührstäbchen, \SI{10}{\milli\liter} vorbereitete Nährstofflösung und 2 Dosiereinheiten Nitrifikationshemmen gegeben bevor die eigentliche Waserprobe hinzugegeben wurde. Der eingelegte Dichtkörper wurde in jeder Flasche mit je 2 Plätzchen Natriumhydroxid bestückt. Die Dichtung des Deckels muss mit etwas Schlifffett bestrichen werden bevor er handfest aufgeschraubt wird. Anschließend werden die Analyseflaschen in den Messschrank auf Magnetrührplatten gestellt und an die Innenliegenden Schläuche angeschlossen. Der Name des genutzten manometrischen Messsystems der Firma Nanotec lautet BSB-Win2000. Der BSB wird über einen Zeitraum von 5 Tagen mittels der Differenzdruckmethode bestimmt.



\chapter{Ergebnisse}
\label{sec:ergebnisse}

Im folgenden Protokollabschnitt werden die Versuchsergebnisse der Versuchsdurchführung präsentiert.




\vspace*{-3.5mm}

\section{Sedimentationsverhalten}
Im Folgenden sind die Messwerte und die dazugehörigen Graphen zum Sedimentationsverhalten in Tabelle \ref{tab:sed.data} und in den Abbildungen \ref{dia:sedimentation_1} und \ref{dia:sedimentation_12} dargestellt. Das Absetzvolumen nach \SI{30}{\minute} ist in den Graphen mittels einer waagerecht verlaufenden, gestrichelten Linie angedeutet.

\vspace*{-2.5mm}
\renewcommand{\arraystretch}{1.2}
% Table generated by Excel2LaTeX from sheet 'Tabelle1'
\begin{table}[h!]
	\centering
	\caption{Sedimentation Messwerte}
	\begin{tabular}{c|ccc}
		\toprule
		\textbf{Zeit [\si{\second}]} & \multicolumn{1}{l}{\textbf{Probe 1 [\si{\milli\liter}]}} & \multicolumn{1}{l}{\textbf{Probe 2 [\si{\milli\liter}]}} & \multicolumn{1}{l}{\textbf{Probe 3 [\si{\milli\liter}]}} \\
		\midrule
		\rowcolor[rgb]{ .851,  .851,  .851} 0     & 1000  & 0,00  & 0,00 \\
		1     & 800   & 0,10  & 0,10 \\
		\rowcolor[rgb]{ .851,  .851,  .851} 2     & 750   & 0,25  & 0,20 \\
		3     & 700   & 0,30  & 0,30 \\
		\rowcolor[rgb]{ .851,  .851,  .851} 4     & 550   & 0,40  & 0,40 \\
		5     & 450   & 0,50  & 0,50 \\
		\rowcolor[rgb]{ .851,  .851,  .851} 6     & 400   & 0,50  & 0,50 \\
		7     & 350   & 0,50  & 1,00 \\
		\rowcolor[rgb]{ .851,  .851,  .851} 8     & 350   & 0,50  & 1,70 \\
		9     & 300   & 0,60  & 2,00 \\
		\rowcolor[rgb]{ .851,  .851,  .851} 10    & 300   & 1,00  & 2,00 \\
		15    & 250   & 1,50  & 3,00 \\
		\rowcolor[rgb]{ .851,  .851,  .851} 20    & 250   & 1,50  & 3,50 \\
		25    & 225   & 1,90  & 4,00 \\
		\rowcolor[rgb]{ .851,  .851,  .851} 30    & 210   & 2,00  & 4,00 \\
		\bottomrule
	\end{tabular}%
	\label{tab:sed.data}%
\end{table}%
\FloatBarrier


% pgf diagramm3 START
\begin{figure}[h!]
	\begin{center}
		%\resizebox{0.8\textwidth}{!}{
		\begin{tikzpicture}[trim axis left, trim axis right]
		\begin{axis}[
		axis lines = left,
		width = 13cm,
		height = 7cm,
		xmin = 0,
		xmax = 35,
		ymin = 0,
		ymax = 1200,
		ylabel={Absetzvolumen in $\si{\milli \liter}$},
		y label style={at={(-0.03,0.5)}},
		xlabel={Zeit in min},
		]
		\addplot[black,mark=*,text mark as node=true,point meta=explicit symbolic,nodes near coords]
		coordinates {(0, 1000) (1, 800) (2, 750) (3, 700) (4, 550) (5, 450) (6,400) (7,350) (8,350) (9,300) (10,300) (15,250) (20,250) (25,225)(30,210)};
		
		\addplot[no markers, dotted] coordinates {(0,210)(35,210)};
		
		\legend{Probe 1, Absetzvolumen (nach 30 min)}
		\end{axis}
		\end{tikzpicture}
		%	}
		\caption{Sedimentationskurve für Abwasserprobe 1}
		\label{dia:sedimentation_1}
	\end{center}
\end{figure}
\FloatBarrier
% pgf diagramm3 ENDE

% pgf diagramm3 START
\begin{figure}[h!]
	\begin{center}
		%\resizebox{0.8\textwidth}{!}{
		\begin{tikzpicture}[trim axis left, trim axis right]
		\begin{axis}[
		axis lines = left,
		width = 13cm,
		height = 7cm,
		xmin = 0,
		xmax = 35,
		ymin = 0,
		ymax = 5.5,
		ytick = {0,0.5,...,5},
		xtick = {0,5,...,30},
		ylabel={Absetzvolumen in $\si{\milli \liter}$},
	%	y label style={at={(-0.03,0.5)}},
		xlabel={Zeit in min},
		legend style={at={(0.5,0.95)},anchor=west}
		]
		
		\addplot[black,mark=*,text mark as node=true,point meta=explicit symbolic,nodes near coords, dashed]
		coordinates {(0, 0) (1, 0.1) (2, 0.25) (3, 0.3) (4, 0.4) (5, 0.5) (6,.5) (7,.5) (8,.5) (9,.6) (10,1) (15,1.5) (20,1.5) (25,1.9)(30,2)};
		
		\addplot[black,mark=x,text mark as node=true,point meta=explicit symbolic,nodes near coords]
		coordinates {(0, 0) (1, 0.1) (2, 0.2) (3, 0.3) (4, 0.4) (5, .5) (6,.5) (7,1.0) (8,1.7) (9,2) (10,2.0) (15,3) (20,3.5) (25,4.)(30,4)};
		
		\addplot[no markers, dotted] coordinates {(0,4)(35,4)};
		
		\addplot[no markers, dotted] coordinates {(0,2)(35,2)};
		
		\legend{Probe 2, Probe 3, Absetzvolumen (nach 30 min)}
		\end{axis}
		\end{tikzpicture}
		%	}
		\caption{Sedimentationskurve für Abwasserproben 2 und 3}
		\label{dia:sedimentation_12}
	\end{center}
\end{figure}
\FloatBarrier
% pgf diagramm3 ENDE



\section*{Absetzvolumen}

%Tabelle START
\vspace*{-2.5mm}
\renewcommand{\arraystretch}{1.2}
\begin{table}[h!]
	%	\begin{adjustwidth}{-2cm}{-2cm} 
	%	\begin{center}
	\centering
	\caption{Messwerte für abfiltrierbare Stoffe}
	\label{tab:absetzvol}
	%	\resizebox{19cm}{!}{
	\begin{tabulary}{1.7\textwidth}{L|CCC}
		\hline
		& \textbf{Probe 1 $\boldsymbol{\left[\si{\milli \liter \per \liter}\right]}$}  & \textbf{Probe 2 $\boldsymbol{\left[\si{\milli \liter \per \liter}\right]}$} & \textbf{Probe 3 $\boldsymbol{\left[\si{\milli \liter \per \liter}\right]}$}  \\ 
		\hline
		\textbf{Absetzvolumen (nach 30 min)}&  210 	& 2 	& 4\\
		\hline
	\end{tabulary}
	%	}
	%	\end{center}
	%	\end{adjustwidth}
\end{table}
\FloatBarrier
\vspace*{-2.5mm}

%Tabelle Ende
\newpage
\section{Trockensubstanz TS}
Beispielrechnung für Probe 1
\begin{flalign}
	m_{\text{Filterkuchen}}		&= m_{\text{Filterkuchen+Filter}}-m_{\text{Filter (trocken)}}\\
	m_{\text{Filterkuchen,1}}	&= \SI{0,611}{\gram}-\SI{0,487}{\gram}\\
								&=\underline{ \SI{0,124}{\gram}}\\[8pt]
	c_{\text{filtrierbar}} 		&= \frac{m_{\text{Filterkuchen}}}{V_{\text{Probe}}}\\[2mm]
	c_{\text{filtrierbar,1}} 	&= \frac{m_{\text{Filterkuchen,1}}}{V_{\text{Probe 1}}}\\[2mm]
								&= \frac{\SI{0,124}{\gram}}{\SI{94}{\milli \liter}}\\
								&= \SI{0.0013191489}{\gram\per \milli \liter} \approx \underline{\underline{\SI{1,32}{ \gram \per \liter}}}
\end{flalign}

%Tabelle START
\vspace*{-2.5mm}
\renewcommand{\arraystretch}{1.2}
\begin{table}[h!]
%	\begin{adjustwidth}{-2cm}{-2cm} 
	%	\begin{center}
	\centering
	\caption{Messwerte für abfiltrierbare Stoffe}
	\label{tab:filter}
%	\resizebox{19cm}{!}{
	\begin{tabulary}{1.7\textwidth}{L|C|CCC}
		\hline
		 & \textbf{Maßeinheit}&	\textbf{Probe 1} & \textbf{Probe 2} & \textbf{Probe 3}  \\ 
		\hline
	 	\textbf{Probenvolumen}& $\boldsymbol{\left[\si{\milli \liter}\right]}$ & 94 	& 300 	& 100\\
		\textbf{Masse Filter (trocken)} &  $\boldsymbol{\left[\si{\gram}\right]}$ & 0,487 & 0,465 & 0,497\\
		\textbf{Masse Filterkuchen + Filter }& $\boldsymbol{\left[\si{\gram}\right]}$	& 0,611 & 0,528 & 0,506\\
		\hline
		\textbf{Masse Filterkuchen} & $\boldsymbol{\left[\si{\gram}\right]}$& 0,124 &0,063&0,009\\
		\hline
		\textbf{Konzentration abfiltrierbare Stoffe}&$\boldsymbol{\left[\si{\gram \per \liter}\right]}$&1,32&0,21&0,09\\
		\hline
	\end{tabulary}
%	}
	%	\end{center}
%	\end{adjustwidth}
\end{table}
\FloatBarrier
\vspace*{-2.5mm}

%Tabelle Ende

\section*{organische Trockensubstanz oTS}
Beispielrechnung für Probe 1
\begin{flalign}
m_{Filter,Glüh}	&= m_{Filter+Tiegel}-m_{Tiegel}\\
				&= \SI{35,7126}{\gram}-\SI{35,714}{\gram}\\
				&=\underline{ \SI{0,0014}{\gram}}\\[8pt]
m_{\text{oTS}}				&= m_{Filterkuchen}-(m_{\text{Tiegel+Probe}}-m_{\text{Tiegel}})-m_{Filter,Glüh}\\
m_{\text{oTS,1}}			&= \SI{0,124}{\gram}-(\SI{33,2209}{\gram}-\SI{33,2023}{\gram})-\SI{0,0014}{\gram}\\
							&= \underline{ \SI{0,104}{\gram}}\\[8pt]
c_{\text{oTS}} 				&= \frac{m_{\text{oTS}}}{V_{\text{Probe}}}\\[2mm]
c_{\text{oTS,1}} 			&= \frac{\SI{0,104}{\gram}}{\SI{94}{\milli \liter}}\\
							&= \SI{0.001106}{\gram\per \milli \liter} \approx \underline{\underline{\SI{1106}{\milli \gram \per \liter}}}
\end{flalign}

\begin{flalign}
		w_{oTS}&=\frac{m_{oTS}}{m_{Filterkuchen}}\\
		w_{oTS,1}&=\frac{m_{oTS,1}}{m_{Filterkuchen}}\\
		&= \frac{\SI{0,104}{\gram}}{\SI{0,124}{\gram}}\\
		&\approx\underline{\underline{ \SI{84}{\percent}}}
\end{flalign}

%Tabelle START
\vspace*{-2.5mm}
\renewcommand{\arraystretch}{1.2}
\begin{table}[h!]
	%	\begin{adjustwidth}{-2cm}{-2cm} 
	%	\begin{center}
	\centering
	\caption{Messwerte für organische Trockensubstanz}
	\label{tab:ots}
	%	\resizebox{19cm}{!}{
	\begin{tabulary}{1.7\textwidth}{L|C|CCC}
		\hline
		& \textbf{Maßeinheit}&	\textbf{Probe 1} & \textbf{Probe 2} & \textbf{Probe 3}  \\ 
		\hline
		\textbf{Probenvolumen}& $\boldsymbol{\left[\si{\milli \liter}\right]}$ & 94 	& 300 	& 100\\
		\textbf{Masser Filter (verglüht)} & $\boldsymbol{\left[\si{\milli \gram}\right]}$ & \SI{0,0014}{\gram}&\SI{0,0014}{\gram}&\SI{0,0014}{\gram}\\
		\textbf{Masse Tiegel (leer)} &  $\boldsymbol{\left[\si{\gram}\right]}$ & 33,2023 & 31,0858 & 34,0671\\
		\textbf{Masse Tiegel (voll)}& $\boldsymbol{\left[\si{\gram}\right]}$	& 33,2209 & 31,1029 & 34,0713\\
		\hline
		\textbf{Masse oTS} & $\boldsymbol{\left[\si{ \gram}\right]}$& 0,104 & 0,106 & 0,118\\
		\hline
		\textbf{Konzentration oTS}&$\boldsymbol{\left[\si{\milli \gram \per \liter}\right]}$& 1106 & 352 & 1184 \\
		\textbf{Massenanteil oTS an TS} & $\boldsymbol{\left[\%\right]}$&84&85&95\\
		\hline
	\end{tabulary}
	%	}
	%	\end{center}
	%	\end{adjustwidth}
\end{table}
\FloatBarrier
\vspace*{-2.5mm}
%Tabelle Ende

\section{Schlammvolumenindex SVI}
%Quelle Römpp \cite{Dr.ManfredNeupert.August2008}\\

Beispielrechnung für Probe 1
\begin{flalign}
	SVI		&= \frac{\varphi_{\text{Absetzvolumen nach 30 min}}}{c_{\text{filtrierbar}}}\\
	SVI_1	&= \frac{\varphi_{\text{Absetzvolumen nach 30 min},1}}{c_{\text{filtrierbar},1}}\\
			&= \frac{\SI{210}{\milli \liter \per \liter}}{\SI{1319}{\milli \gram \per \liter}}\\
			&= \SI{0.1592}{\milli \liter \per \milli \gram} = \underline{\underline{\SI{159,2}{\milli \liter \per \gram}}}
\end{flalign}
%Tabelle START
\vspace*{-2.5mm}
\renewcommand{\arraystretch}{1.2}
\begin{table}[h!]
	%	\begin{adjustwidth}{-2cm}{-2cm} 
	%	\begin{center}
	\centering
	\caption{SVI für die Abwasserproben 1 bis 3}
	\label{tab:svi}
	%	\resizebox{19cm}{!}{
	\begin{tabulary}{1.7\textwidth}{L|C|C|C}
		\hline
		& \textbf{Probe 1 $\boldsymbol{\left[\si{\milli \liter \per \gram}\right]}$}& \textbf{Probe 2 $\boldsymbol{\left[\si{\milli \liter \per \gram}\right]}$} & \textbf{Probe 3 $\boldsymbol{\left[\si{\milli \liter \per \gram}\right]}$}  \\ 
		\hline
		\textbf{$\boldsymbol{SVI}$}& 159,2 & 9,5 & 44,4 \\
		\hline
	\end{tabulary}
	%	}
	%	\end{center}
	%	\end{adjustwidth}
\end{table}
\FloatBarrier

\newpage

\section{chemischer Sauerstoffbedarf CSB}
Nachfolgend ist die Berechnung des CSB für die Abwasserprobe 1 aufgeführt, da diese als einzige verdünnt analysiert wurde.
Die Verdünnung der Probe mit destilliertem Wasser betrug 1:1. Die CSB-Messampullen wurden für einen Messbereich von 0-1500 mg/l gewählt.
Die Ergebnisse de Analyse sind in Tabelle \ref{tab:csb} dargestellt.
\begin{flalign}
\frac{c_{\text{Probe 1}}}{c_{\text{Probe 1 (verdünnt)}}}  	&= \frac{2}{1} \\
c_{\text{Probe 1}}											&= 2*c_{\text{Probe 1 (verdünnt)}}\\
															&= 2*\SI{575}{\milli \gram \per \liter} \\
															&= \underline{\underline{\SI{1150}{\milli \gram \per \liter}}}
\end{flalign}
%Tabelle START
\vspace*{-2.5mm}
\renewcommand{\arraystretch}{1.2}
\begin{table}[h!]
	%	\begin{adjustwidth}{-2cm}{-2cm} 
	%	\begin{center}
	\centering
	\caption{Messwerte für den chemischen Sauerstoffbedarf für die Abwasserproben 1 bis 3}
	\label{tab:csb}
	%	\resizebox{19cm}{!}{
	\begin{tabulary}{1.7\textwidth}{L|CC|C|C}
		\hline
		& \textbf{Probe 1 (verdünnt) $\boldsymbol{\left[\si{\milli \gram \per \liter}\right]}$}  &\textbf{Probe 1 $\boldsymbol{\left[\si{\milli \gram \per \liter}\right]}$}& \textbf{Probe 2 $\boldsymbol{\left[\si{\milli \gram \per \liter}\right]}$} & \textbf{Probe 3 $\boldsymbol{\left[\si{\milli \gram \per \liter}\right]}$}  \\ 
		\hline
		\textbf{$\boldsymbol{CSB}$}& 575 & 1150 & 486 & 221 \\
		\hline
	\end{tabulary}
	%	}
	%	\end{center}
	%	\end{adjustwidth}
\end{table}
\FloatBarrier

\vspace*{-2.5mm}

\section[biologischer Sauerstoffbedarf BSB$_5$]{biologischer Sauerstoffbedarf BSB$\boldsymbol{_5}$}
Als Messberereich für Probe 1 und 2 wurde 0-\SI{700}{\milli \gram\per\liter} gewählt. Das Probenvolumen lag jeweils bei \SI{95}{\milli\liter}. Für Probe 3 wurde aufgrund des niedriger zu erwartenden $BSB_5$ der nächst geringere Messbereich von 0-\SI{350}{\milli\gram\per\liter} verwendet. Das Probenvolumen lag für Probe 3 bei \SI{160}{\milli\liter}.
Die für die Messbereiche geforderten Volumina sind der Tabelle unter \cite[S.22]{Skript} entnommen worden.\\
Die Messwerte zum $BSB_5$ für die Abwasserproben 1 und 3 sind in Tabelle \ref{tab:bsb} aufgeführt.

%Tabelle START
\vspace*{-2.5mm}
\renewcommand{\arraystretch}{1.2}
\begin{table}[h!]
	%	\begin{adjustwidth}{-2cm}{-2cm} 
	%	\begin{center}
	\centering
	\caption{Messwerte für den biologischen Sauerstoffbedarf über 5 Tage für die Abwasserproben 1 bis 3}
	\label{tab:bsb}
	%	\resizebox{19cm}{!}{
	\begin{tabulary}{1.7\textwidth}{L|C|C|C}
		\hline
		& \textbf{Probe 1 $\boldsymbol{\left[\si{\milli \gram \per \liter}\right]}$}& \textbf{Probe 2 $\boldsymbol{\left[\si{\milli \gram \per \liter}\right]}$} & \textbf{Probe 3 $\boldsymbol{\left[\si{\milli \gram \per \liter}\right]}$}  \\ 
		\hline
		\textbf{$\boldsymbol{BSB_5}$}& 254 & 252 & 50 \\
		\hline
	\end{tabulary}
	%	}
	%	\end{center}
	%	\end{adjustwidth}
\end{table}
\FloatBarrier


\newpage

\section{Gegenüberstellung der Mindestanforderungen für das Einleiten kommunaler Abwässer in einen Vorfluter der GK 5 mit den Abwasserproben}
Die Referenzwerte der Mindestanforderungen für das Einleiten kommunaler Abwässer in den Vorfluter der Größenklasse 5 sind im Anhang von \cite[S. 29]{Skript} zu finden.
\vspace*{-2.5mm}
\renewcommand{\arraystretch}{1.2}
\begin{table}[h!]
	\centering
	\caption{Tabellarischer Vergleich der Messwerte mit den Mindestanforderungen für das Einleiten kommunaler Abwässer in den Vorfluter der GK 5}
	\label{tab_vgl}
	%\resizebox{10cm}{!}{
	\begin{tabulary}{1.2\textwidth}{l|C|C}
		\hline
		 & \textbf{$CSB$} $\boldsymbol{\left[\si{\milli\gram\per\liter}\right]}$ & \textbf{$BSB_5$} $\boldsymbol{\left[\si{\milli\gram\per\liter}\right]}$\\
		\hline
		\textbf{Grenzwert} & \textbf{75} & \textbf{15}  \\
		\hline
		Probe 1 & 575 & 254 \\
		Probe 2 & 486 & 252\\
		Probe 3 & 221 & 50\\
		\hline
	\end{tabulary}
	%}
\end{table}
\FloatBarrier

\begin{figure}[h!]
	\begin{tikzpicture}
	\selectcolormodel{gray}
	\begin{axis}[
	xbar=1pt,% space of 0pt between adjacent bars
	bar width=7,
	width=15cm,
	height=7cm,
	%minor y tick num=4,
	xmax=600,xmin=0,
	x tick label style={/pgf/number format/.cd,%
		scaled x ticks = false,
		set decimal separator={,},
		fixed},
	symbolic y coords={Probe 3,Probe 2,Probe 1,max. für GK 5},
	ytick=data,
	%legend style={at={(0.5,-0.15)},
	%	anchor=north,legend columns=-1},
	xtick={0,50,...,600},
	grid=major,
	xlabel=Gehalt in  \si{\milli\gram\per\liter},
	%enlargelimits=0.15,
	%postaction={pattern=north east lines}
	]
	%CSB
	\addplot[fill=black] coordinates {
		(575,Probe 1) (486,Probe 2) (221,Probe 3) (75,max. für GK 5)
	};
	%BSB5
	\addplot coordinates {
		(254,Probe 1) (252,Probe 2) (50,Probe 3) (15,max. für GK 5)
	};

	\legend{$CSB$,$BSB_5$}
	\end{axis}
	\end{tikzpicture}
	\caption{Vergleich mit Mindestanforderungen für das Einleiten kommunaler Abwässer in den Vorfluter der GK 5 für die Abwasserproben 1 bis 3}
	\label{Balkendiagramm}
\end{figure}
\FloatBarrier


\newpage

\section{Gegenüberstellung der durchschnittlichen Beschaffenheit von häuslichem Abwasser mit den Abwasserproben}

Um die Messwerte des Versuches mit häuslichem Abwasser gegenüberzustellen wird die Tabelle Tab. \ref{tab:komm} (siehe \cite[S. 29]{Skript}) genutzt.

\vspace*{0.5cm}
\renewcommand{\arraystretch}{1.2}
\begin{table}[h!]
	\centering
	\caption[Tabellenausschnitt zur durchschnittlichen Beschaffenheit von häuslichem Abwasser]{Tabellenausschnitt zur durchschnittlichen Beschaffenheit von häuslichem Abwasser \cite[S. 29]{Skript}}
	\label{tab:komm}
	%\resizebox{10cm}{!}{
	\begin{tabulary}{1.2\textwidth}{l|C|C|C|C}
	\textbf{Kriterium} 		& \textbf{Maßeinheit} 				&\multicolumn{3}{c}{\textbf{Belastungsgrad}}\\
	\hline
							&									& gering	& mittel & stark\\
	\hline
	Absetzbare Stoffe		&\si{\milli \liter \per \liter} 	& 2			& 6	 	 & 12\\
	Abfiltrierbare Stoffe	&\si{\milli \gram \per \liter} 		& 200		&500	 & 900\\
	$CSB$					&\si{\milli \gram \per \liter} 		& 300		&600	 & 1000\\
	$BSB_5$					&\si{\milli \gram \per \liter}		& 150		&300	 & 500\\
	\end{tabulary}
	%}
\end{table}
\FloatBarrier
\vspace*{1.5cm}

\begin{figure}[h!]
	\begin{tikzpicture}
	\begin{axis}[
	%x tick label style={
	%	/pgf/number format/1000 sep=},
	xtick = data,
	ylabel=Absetzbare Stoffe in \si{\milli \liter \per \liter},
   	enlarge x limits=0.15,
	ybar,
	ymin = 0,
	ymax = 20,
	bar width=10pt,
	width=15cm,
	height=9cm,
	symbolic x coords={0,Probe 1, Probe 2,Probe 3,2,1},
	]
	\addplot[fill=gray!50,draw=black!60] coordinates {(Probe 1,20) (Probe 2,2) (Probe 3,4)
	};
	\addplot[green!50!black,sharp plot,update limits=false, dashed] 
	coordinates {(0,2) (1,2)} 
	node[above] at (axis cs:Probe 2,2) {\hspace*{-5.5cm}geringe Belastung
	};
	\addplot[orange,sharp plot,update limits=false, dashed] 
	coordinates {(0,6) (1,6)} 
	node[above] at (axis cs:Probe 2,6) {\hspace*{-5.5cm}mittelere Belastung
	};
	\addplot[red,sharp plot,update limits=false, dashed] 
	coordinates {(0,12) (1,12)} 
	node[above] at (axis cs:Probe 2,12) {\hspace*{-5.5cm}starke Belastung
	};
	\node[red] at (axis cs: Probe 2,18) {\hspace*{-8.cm}\textbf{---\SI{210}{\milli \liter \per \liter}}};
	\legend{Absetzbare Stoffe}
	\end{axis}
	\end{tikzpicture}
	\caption{Absetzbare Stoffe der Abwasserproben 1 bis 3}
	\label{dia:absetz}
\end{figure}
\FloatBarrier

\begin{figure}[h!]
	\begin{tikzpicture}
	\begin{axis}[
	%x tick label style={
	%	/pgf/number format/1000 sep=},
	xtick = data,
	ylabel=Abfiltrierbare Stoffe in \si{\milli \gram \per \liter},
	enlarge x limits=0.15,
	ybar,
	ymin = 0,
	ymax = 1350,
	ytick ={0,200,500,900,1350},
	bar width=10pt,
	width=15cm,
	height=6.7cm,
	symbolic x coords={0,Probe 1, Probe 2,Probe 3,2,1},
	]
	\addplot[fill=gray!50,draw=black!60] coordinates {(Probe 1,1319) (Probe 2,210) (Probe 3,90)
	};
	\addplot[green!50!black,sharp plot,update limits=false, dashed] 
	coordinates {(0,200) (1,200)} 
	node[above] at (axis cs:Probe 2,200) {geringe Belastung
	};
	\addplot[orange,sharp plot,update limits=false, dashed] 
	coordinates {(0,500) (1,500)} 
	node[above] at (axis cs:Probe 2,500) {mittelere Belastung
	};
	\addplot[red,sharp plot,update limits=false, dashed] 
	coordinates {(0,900) (1,900)} 
	node[above] at (axis cs:Probe 2,900) {starke Belastung
	};
	\legend{Abfiltrierbare Stoffe}
	\end{axis}
	\end{tikzpicture}
	\caption{Abfiltrierbare Stoffe der Abwasserproben 1 bis 3 (siehe Tab. \ref{tab:filter})}
	\label{dia:abfilt}
\end{figure}
\FloatBarrier

\begin{figure}[h!]
	\begin{tikzpicture}
	\begin{axis}[
	%x tick label style={
	%	/pgf/number format/1000 sep=},
	xtick = data,
	ylabel=$CSB$ in \si{\milli \gram \per \liter},
	enlarge x limits=0.15,
	ybar,
	ymin = 0,
	ymax = 1400,
	ytick ={0,300,600,1000,1400},
	bar width=10pt,
	width=15cm,
	height=6.7cm,
	symbolic x coords={0,Probe 1, Probe 2,Probe 3,2,1},
	]
	%Daten der Proben
	\addplot[fill=gray!50,draw=black!60] coordinates {(Probe 1,575) (Probe 2,186) (Probe 3,221)
	};
	%Grenzen
	\addplot[green!50!black,sharp plot,update limits=false, dashed] 
	coordinates {(0,300) (1,300)} 
	node[above] at (axis cs:Probe 2,300) {geringe Belastung
	};
	\addplot[orange,sharp plot,update limits=false, dashed] 
	coordinates {(0,600) (1,600)} 
	node[above] at (axis cs:Probe 2,600) {mittelere Belastung
	};
	\addplot[red,sharp plot,update limits=false, dashed] 
	coordinates {(0,1000) (1,1000)} 
	node[above] at (axis cs:Probe 2,1000) {starke Belastung
	};
	\legend{CSB}
	\end{axis}
	\end{tikzpicture}
	\caption{Chemischer Sauerstoffbedarf (CSB) der Abwasserproben 1 bis 3}
	\label{dia:csb}
\end{figure}
\FloatBarrier

\begin{figure}[h!]
	\begin{tikzpicture}
	\begin{axis}[
	%x tick label style={
	%	/pgf/number format/1000 sep=},
	xtick = data,
	ylabel=$BSB_5$ in \si{\milli \gram \per \liter},
	enlarge x limits=0.15,
	ybar,
	ymin = 0,
	ymax = 800,
	ytick ={0,150,300,500,700},
	bar width=10pt,
	width=15cm,
	height=6.7cm,
	symbolic x coords={0,Probe 1, Probe 2,Probe 3,2,1},
	]
	%Daten
	\addplot[fill=gray!50,draw=black!60] coordinates {(Probe 1,254) (Probe 2,252) (Probe 3,50)
	};
	%Linien
	\addplot[green!50!black,sharp plot,update limits=false, dashed] 
	coordinates {(0,150) (1,150)} 
	node[above] at (axis cs:Probe 2,150) {\hspace*{5cm} geringe Belastung
	};
	\addplot[orange,sharp plot,update limits=false, dashed] 
	coordinates {(0,300) (1,300)} 
	node[above] at (axis cs:Probe 2,300) {\hspace*{5cm} mittelere Belastung
	};
	\addplot[red,sharp plot,update limits=false, dashed] 
	coordinates {(0,500) (1,500)} 
	node[above] at (axis cs:Probe 2,500) {\hspace*{5cm} starke Belastung
	};
	\legend{BSB$_5$}
	\end{axis}
	\end{tikzpicture}
	\caption{Biochemischer Sauerstoffbedarf über 5 Tage (BSB$_5$) der Abwasserproben 1 bis 3}
	\label{dia:bsb}
\end{figure}
\FloatBarrier



\chapter{Diskussion}
\label{sec:diskussion}
In diesem Abschnitt des Protokolls werden nun die Ergebnisse des Abschnittes \ref{sec:ergebnisse} diskutiert und ausgewertet.\\\\

%Die erhaltenen Analysenergebnisse sind tabellarisch und grafisch darzustellen.  

%Die Werte sind den Mindestanforderungen für das Einleiten kommunaler Abwässer (GK 5; siehe 
%Anlage des Praktikumsheftes) gegenüberzustellen und zu diskutieren. 
%Ordnen Sie die Proben hinsichtlich ihres Belastungsgrades ein!
  
%Streichen Sie signifikante Unterschiede zwischen den einzelnen Proben heraus und schließen Sie auf 
%ihre Herkunft! 

%Was bedeuten CSB und BSB5 ? Erklären Sie die Merkmale und den Stellenwert beider Analysen.
 
%Warum beträgt die Dauer des BSB-Versuches 5 Tage? 

%Welches CSB/BSB5-Verhältnis besitzt biologisch gut abbaubares Abwasser, welches 
%„Problemwässer“? 

%Errechnen Sie das bestehende CSB/BSB5-Verhältnis und beurteilen Sie die 
%biologische Abbaubarkeit der einzelnen Proben. 

%Geben Sie Empfehlungen zu einer anforderungsgerechten Behandlung der analysierten Proben. 

%Berechnen Sie für die Abwasserschlammproben den Schlammvolumenindex (SVI)! 
%Nach MOHLMANN gibt dieser das Volumen von 1 g Trockensubstanz des Belebtschlammes nach einer 
%Absetzzeit von 30 min an und stellt sich als Verhältnis spez. Volumen des abgesetzten Schlammes in 
%ml/1000ml und spez. Masse der Trockensubstanz in g/1000ml dar. 

%Welche Aussage können Sie mit dem SVI für Ihre Proben machen? 




\chapter{Fehlerbetrachtung}
\label{sec:fehler}
In diesem Abschnitt erfolgt die Fehlerbetrachtung des Versuches, welche Einfluss auf die Messergebnisse haben können.\\

Mess- und Ablesefehler können bei der Bestimmung des Absetzvolumens im Imhofftrichter, bei der automatisierten Messung des CSB, bei der Bestimmung des BSB nach der Differenzdruckmethode, bei allen durchgeführten Wägevorgängen und Volumenmessungen mit Messzylindern und Eppendorfpipetten aufgetreten sein. Auch die Zeitmessung während des Absetzvorgangs ist Fehlerbehaftet, da das Ablesen einige Zeit in Anspruch nahm. Darum handelt es sich bei allen Zeitangaben nur um Näherungswerte mit Abweichungen von bis zu 30 Sekunden.\\
Die Belastungsgrade, welchen die Proben zugeordnet wurden, sind nicht klar definiert. Es handelt sich stets um interpretationsbedürftige Bereiche. So verhält es sich auch mit den Referenzwerten des $\frac{CSB}{BSB_5}$-Verhältnisses. Je nach Quelle variieren diese geringfügig. \\
Die Versuchsdurchführung weist an einigen Stellen ein beachtenswertes Fehlerpotential auf. Die Trocknung des Filterpapiers auf der Trocknungswage ist recht genau, da die Wage den Gewichtsverlust kontrolliert und erst bei Stagnation der Masse die Trockenmasse anzeigt. Die Vorbereitung der CSB-Ampullen erforderte die Herstellung einer verdünnten Probenlösung. Hierbei kamen nur Messzylinder zum Einsatz. Dem entsprechend kann das angegebene Volumen um wenige Milliliter variieren. Die eigentliche Befüllung der Ampullen mit Eppendorfpipetten dürfte keinen signifikanten Fehler verursacht haben. \\ 
Beim Verglühen der Filterkuchen wurden die Filterpapiere mitverglüht. Um die so zusätzlich eingetragene Masse an Asche herausrechnen zu können, wurde ein Blatt Filterpapier der selben Charge parallel mit verglüht. Ein Mittelwert aus einer größeren Anzahl verglühter Filterpapierblättchen würde einen belastbareren Wert ergeben, da er die natürliche Streuung der Aschegehalte ausgleicht. Das Verglühen der Proben könnte unvollständig gewesen sein. \\
Die Bestimmung des $BSB_5$ ist stark vom Wohlbefinden der Mikroorganismen abhängig. Eine Fehldosierung von Nähstoffen in der Nährlösung, zu hohe Temperaturschwankungen oder das versehentliche Einbringen von Natriumhydroxid könnte die Lebensumstände der Mikroorganismen in einen suboptimalen Bereich verschieben und damit deren Aktivität hemmen. Zur besseren Vergleichbarkeit der $BSB_5$-Werte, hätte die Bakterienpopulation der Proben zu Beginn der Messung aneinander angeglichen werden müssen. Eine niedrige Anzahl von Destruenten im Abwasser verlängert die Abbaudauer dadurch, dass sich die Population erst entwickeln muss. Im ungünstigsten Fall reicht dann die Zeit von fünf Tagen nicht zum weitest möglichen biologischen Abbau aus. Undichtigkeiten im System würden die Differenzdruckmessung kompromittieren. 


%Praktikumsskript, Modul ………, Versuch …….., Prof. Musterprof. 
%DIN 12345, Jahr der Veröffentlichung 
%Link der Internetseite, Zugriffsdatum 
%Buchtitel, Autor, Verlag, Veröffentlichungsjahr 


%Literaturverzeichnis Bücher
\bibliography{Literatur}
\bibliographystyle{unsrtdin}
\addcontentsline{toc}{chapter}{Literaturverzeichnis}







%Anhang
%\addcontentsline{toc}{chapter}{Anhang}

%\chapter*{Eidesstattliche Erklärung}
\label{erklaerung}
Hiermit versichere ich, die vorliegende Seminararbeit selbstständig und nur unter Verwendung der von mir angegebenen Quellen und Hilfsmittel verfasst zu haben. Sowohl inhaltlich als auch wörtlich entnommene Inhalte wurden als solche kenntlich gemacht. Die Arbeit hat in dieser oder vergleichbarer Form noch keinem anderem Prüfungsgremium vorgelegen. \\
\\[1.5cm]
Datum:	\hrulefill\enspace Unterschrift: \hrulefill
\\[3.5cm]
\addcontentsline{toc}{chapter}{Selbstständigkeitserklärung}

%\chapter{Bausteine}

\section{Beispiel für Tabelle}

%Tabelle START
\vspace*{-2.5mm}
\renewcommand{\arraystretch}{1.2}
\begin{table}[h!]
	\centering
	\caption{Abmessungen der Probekörper vor dem Zugversuch}
	\label{tab:tabelle1}
	%\resizebox{10cm}{!}{
	\begin{tabulary}{\textwidth}{C|CCC}
		\hline
		\textbf{Probe}  &\textbf{Breite [mm]}&\textbf{Dicke [mm]}&\textbf{Anf.-länge[mm]} \\ 
		\hline
		Kupfer (gewalzt) & 12,5 &3,00&50\\
		Kupfer (geglüht) & 12,5&3,00&50\\
		PA6 & 10,0&4,00&50\\
		PP (EPR-30\% Kautschuk) & 9,9 &3,95&50\\
		\hline
	\end{tabulary}
	%}
\end{table}

\FloatBarrier
\vspace*{-2.5mm}
%Tabelle ENDE

\section*{Tabelle mit Itemize}
%Tabelle START
\vspace*{-2.5mm}
\renewcommand{\arraystretch}{1.2}
\begin{table}[h!]
	\centering
	\caption*{Vor- und Nachteile der Geothermie}
	\label{tab:tabelle1}
	\begin{tabulary}{\textwidth}{C|C}
		\hline
		\textbf{Vorteile}  &\textbf{Nachteile} \\ 
		\hline
		&\\
		\begin{minipage}[t]{0.4\textwidth}
			\begin{itemize}
				\item Strom, Wärme und Kälte wird erzeugt
				\item keine saisonalen und tageszeitlichen Schwankungen
				\item 	quasi-regenerativ
				\item 	nachfrage-gerechte Energiebereitstellung
				\item Erzeugungspotenzial sehr hoch 
				\item 	grundsätzlich standortunabhängig
			\end{itemize}
		\end{minipage} & 
		\begin{minipage}[t]{0.4\textwidth}
			\begin{itemize}
				\item hohe Anschaffungskosten
				\item abhängig von geologischen Gegebenheiten
				\item geringer Stromwirkungsgrad (thermodynamisch bedingt)
				\item keine Marktdurchdringung in DE
				\item erfahrene Bauunternehmen notwendig
				\item gute Vorerkundung und Überwachung notwendig
			\end{itemize}
		\end{minipage}\\
	\end{tabulary}
\end{table}
\FloatBarrier
\vspace*{-2.5mm}
%Tabelle ENDE

\newpage

\section*{Beispiel für Skalierbare Tabelle}
%TAbelle Start
\vspace*{-2.5mm}
\renewcommand{\arraystretch}{1.2}
\begin{table}[h!]
	\centering
	\caption*{}
	\resizebox{0.5\textwidth}{!}{
		\begin{tabulary}{\textwidth}{C|C|C|C}
			\textbf{Name} & \textbf{Anwendung}&\textbf{Gleichung}&\textbf{Stoffkonstante} \\ 
			\hline  
			KICK& $x_{80_\omega}>\SI{50}{\milli\meter}$ &$e_{KICK}=c_K*log(\frac{x_{80_\omega}}{x_{80_\alpha}})$&$c_K=1,15*\frac{c_B}{\sqrt{0,05\si{\meter}}} \left[\si{\raiseto{2}\meter\per\raiseto{2}\second}\right]$\\
			BOND&$\SI{50}{\micro\meter}<x_{80_\omega}<\SI{50}{\milli\meter}$&$e_{BOND}=c_B*\left(\frac{1}{\sqrt{x_{80_\omega}}}-\frac{1}{\sqrt{x_{80_\alpha}}}\right)$& $c_B$: tabelliert $\left[\si{\raiseto{2,5}\meter\per\raiseto{2}\second}\right]$\\ 
			RITTER& $x_{80_\omega}>\SI{50}{\micro\meter}$&$e_{RITT}=c_R*\left(\frac{1}{x_{80_\omega}}-\frac{1}{x_{80_\alpha}}\right)$&$c_R= 0,5*c_B*\sqrt{\SI{5e-5}{\meter}}$ \\  
	\end{tabulary}}
\end{table}
\FloatBarrier
%Ende TAbelle

\section{Befehlszeilen in Text einfügen}
Befehlzeilen aus der "Latex-Sprache"\ lassen sich nicht ohne weiteres im Text darstellen. Das System erkennt diese als solche und gibt warnhinweise aus. In der Regel werden die Befehle dann auch falsch dargestellt. Umgehen lässt sich diese Problematik mit der verbatim-Umgebung. In ihren Grenzen werden Eingaben 1:1 dargestellt wie eingegeben. Die Funktionsweise soll nachfolgend durch ein Beispiel verdeutlicht werden.
\begin{verbatim*}

\begin{verbatim}
\usepackage{Beispieltrolllllollllolll} 
\end{verbatim}

\end{verbatim*}

Sollte besonderer Wert auf die Kenntlichmachung der Leerzeichen gelegt werden kann auch mit \texttt{verbatim*} gearbeitet werden.

\section{Seitenübergreifende, lange Tabellen}

Tabellen welche Messwerte in einem solchen Umfang enthalten, dass sie nicht auf einer einzelnen A4-Seite Platz finden, können mit dem Paket \begin{verbatim}
\usepackage{longtable} 
\end{verbatim}
in ein Dokument eingepflegt werden wie folgendes Beispiel belegt.:
\begin{longtable}[c]{lllll}
	\caption{Dehnungstabelle}\\
	\label{alles}
	$Zeit [HH:MM:SS]$ & $\Delta l_{PE} [mm]$ & $\varepsilon_{PE}$ & $\Delta l_{Pb} [mm]$ & $\varepsilon_{Pb}$ \\
	\hline
	\endfirsthead
	%\caption{}\\
	$Zeit [HH:MM:SS]$ & $\Delta l_{PE} [mm]$ & $\varepsilon_{PE}$ & $\Delta l_{Pb} [mm]$ & $\varepsilon_{Pb}$ \\ 
	\hline
	\endhead
	\multicolumn{5}{r}{Fortsetzung auf n{\"a}chster Seite}\\
	\endfoot
	\hline
	\multicolumn{5}{r}{} \\
	\endlastfoot
	% Ab hier kommt der Inhalt der Tabelle
	00:00:00 & 3,00 & 0,11 & 3,30 & 0,12\\
	00:00:10 & 3,80 & 0,14 & 3,42 & 0,13\\
	00:00:20 & 4,15 & 0,16 & 3,53 & 0,13\\
	00:00:30 & 4,43 & 0,17 & 3,68 & 0,14\\
	00:00:40 & 4,61 & 0,17 & 3,71 & 0,14\\
	00:00:50 & 4,80 & 0,18 & 3,81 & 0,14\\
	00:01:00 & 4,95 & 0,19 & 4,12 & 0,15\\
	00:20:00 & 8,03 & 0,30 &  & \\
\end{longtable} 

\section{Diagramme}
Diagramme lassen sich unter anderem mit dem Paket 
\begin{verbatim}
\usepackage{pgfplots} 
\end{verbatim}
implementieren.
Um das Dokument einigermaßen klein zu halten empfehle ich folgendes weiters Paket zu nutzen.
\begin{verbatim}
\usepackage{csvsimple} 
\end{verbatim}
Es erleichter das importieren von Datensätzen aus .csv Dateien. Die .csv Datei muss folgenden Anforderungen genügen:- Dezimaltrenner Punkt
\begin{itemize}
	\item Dezimaltrenner Punkt
	\item Jede Zeile umgebrochen
	\item Koordinaten durch komma getrennt
	\item Spalten beschriften z.B x,y oder a,b
\end{itemize}

Sollte man eine Tabellenkalkulationsdatei in eine csv Datei konvertiert haben, hat dieses meist die Falsche Form. Durch die Funktion suchen und ersetzen (bsp emacs) können aber sehr schnell die nötigen Korrekturen erfolgen. 

\begin{figure}[h]
	\begin{center}
		\begin{tikzpicture}
		\begin{axis}[
		width=12cm,
		height=6cm,
		xlabel=Zeit in Sekunden,
		ylabel=Dehnung]
		
		\begin{scope}[brown]
		\draw[brown] ({axis cs:10,0}|-{rel axis cs:0,1}) -- ({axis cs:10,0}|-{rel axis cs:0,0});
		\draw[brown] ({axis cs:120,0}|-{rel axis cs:0,1}) -- ({axis cs:120,0}|-{rel axis cs:0,0});
		\end{scope} 
		
		\addplot table [x=a, y=b, col sep=comma] {data/KriechkurvePb.csv};
		\end{axis}
		\end{tikzpicture}
		\caption{Kriechkurve Blei}
		\label{kkb}
	\end{center}
\end{figure} 
\FloatBarrier                     

\section*{Beispiel für Berechnungen}
Die Berechnung der wahren Spannung $\sigma_{W}$ bei Höchstkraft erfolgt unter der Annahme, dass der Prüfkörperquerschnitt noch 60\% des Ausgangsquerschnitts beträgt. Die Berechnung erfolgt ab Gleichung \ref{ber1}.\\ 

%Berechnung der Fläche A
\textbf{Fläche} $\boldsymbol{A}$ \textbf{:}
%Start
\begin{flalign}
A 	&= \frac{\pi}{4}*d^2\\
&=\frac{\pi}{4}*(\SI{80}{\milli \meter})^2
\end{flalign}
%Ende

\section{Beispiel für ein Bild}

%Start
\begin{figure}[h!]
	\centering
	\includegraphics[width=0.60\textwidth]{img/skizzepruef3}
	\caption{Skizze Prüfkörperbemaßung}
	\label{skizzepruef}
\end{figure}
\FloatBarrier
%Ende

\newpage

\section{Beispiel für zwei Bilder}
\label{sec:versuchsaufbau}

%Start
\begin{figure}[h!]
	\centering
	\begin{subfigure}{.5\textwidth}
		\centering
		\includegraphics[width=0.75\textwidth]{Aufbau2}
		\caption{Skizze zum Versuchsaufbau}
		\label{fig:sub1}
	\end{subfigure}%
	\begin{subfigure}{.5\textwidth}
		\centering
		\includegraphics[width=0.6\textwidth]{img/Aufbau1}
		\caption{realer Versuchsaufbau}
		\label{fig:sub2}
	\end{subfigure}
	\caption{Versuchsaufbau als Skizze und in Realität}
	\label{fig:aufbau} 
\end{figure}
\FloatBarrier
%Ende

\newpage


\section{Beispiel für vier Bilder}
%Start
\begin{figure}[h!]
	\centering
	\begin{subfigure}{.5\textwidth}
		\centering
		\includegraphics[width=0.75\textwidth]{img/Kupfer_gewalzt}
		\caption{Kupfer (gewalzt)}
		\label{fig:sub3}
	\end{subfigure}%
	\begin{subfigure}{.5\textwidth}
		\centering
		\includegraphics[width=0.75\textwidth]{img/Kupfer_weich}
		\caption{Kupfer (geglüht)}
		\label{fig:sub4}
	\end{subfigure}
	\begin{subfigure}{.5\textwidth}
		\centering
		\includegraphics[width=0.75\textwidth]{img/PA6}
		\caption{PA6}
		\label{fig:sub5}
	\end{subfigure}%
	\begin{subfigure}{.5\textwidth}
		\centering
		\includegraphics[width=0.75\textwidth]{img/PP}
		\caption{PP (ERP-30\% Kautschuk)}
		\label{fig:sub6}
	\end{subfigure}
	
	\caption{Bruchstellennahaufnahmen der Probekörper}
	\label{fig:bruchstellen} 
\end{figure}
\FloatBarrier
%Ende

\section{Beispiel Einheiten}

\begin{align*}
\textbf{$\SI{12,0/12}{\kg\meter\per\second \raiseto{5} \per \xyz}*\SI{13}{\per\raiseto{-2}\meter}=\SI{256}{}$}
\end{align*}

\begin{align}
\SI{12,0/12}{\meter\per\joule}*\SI{13}{\gram}=\SI{256}{\coulomb}
\end{align}

\newpage

\section{Beispiel für Mini-Formelsammlung}
\begin{flalign}
\label{gl1}
\text{\textbf{Dehnung (Def.)} } \boldsymbol{\varepsilon} \text{ \textbf{:}} && \hspace*{-1em}  \varepsilon=\frac{\Delta l}{l_0} &&
\end{flalign}

\begin{flalign}
\label{gl2}
\text{\textbf{norminelle Spannung} } \boldsymbol{\sigma} \text{\textbf{:}} && \hspace*{-3em} \sigma=\frac{F}{A_0} &&
\end{flalign}

\begin{flalign}
\label{gl3}
\text{\textbf{Sekantenmodul (Kunststoffe) }} \boldsymbol{E_S} \text{ \textbf{:}} && E_S=\frac{\sigma_2-\sigma_1}{\varepsilon_2-\varepsilon_1}=\frac{F_2-F_1}{0,002*A_0} &&
\end{flalign}

\begin{flalign}
\label{gl4}
\text{\textbf{E-Modul (Metalle) }} \boldsymbol{E_M} \text{ \textbf{:}} && \hspace*{5em} E_M=\frac{\sigma}{\varepsilon}=\frac{\sigma_2-\sigma_1}{\varepsilon_2-\varepsilon_1}=\frac{R_{p_{0.2\%}}}{0,2\%} &&
\end{flalign}

\begin{flalign}
\label{gl5}
\text{\textbf{Bruchdehnung} } \boldsymbol{A} \text{\textbf{:}} && \hspace*{6em} A= \frac{l_{u}-l_{0}}{l_{0}}*100\% &&
\end{flalign}

\begin{flalign}
\label{gl6}
\text{\textbf{Ausgangsquerschnitt } } \boldsymbol{S_0} \text{\textbf{:}} && \hspace*{3em} S_{0}= Breite*Dicke &&
\end{flalign}
\begin{flalign}
\label{gl7}
\text{\textbf{wahre Spannung} } \boldsymbol{\sigma_{W}} \text{\textbf{:}} &&\hspace*{1em} \sigma_{W}=\frac{F_{max}}{S_{End}}&&
\end{flalign}
\begin{flalign}
\label{gl8}
\text{\textbf{Brucheinschnürung } } \boldsymbol{Z} \text{\textbf{:}} && \hspace*{5em} Z=\frac{S_0-S_u}{S_{o}}*100\% &&
\end{flalign}

\newpage

\section*{Fußnoten}
%Start
\begin{figure}[h!]
	\centering
	\includegraphics[width=0.85\textwidth]{tabdia/kfwerte}
	\caption*{$\text{k}_\text{f}$-Werte der Proben 1 bis 11 \protect\footnotemark[1]}
	\label{}
\end{figure}
\FloatBarrier
%Ende

\footnotetext[1]{bezogen auf $V=\SI{50}{\milli \liter}$ und $h=\SI{10}{\milli \meter}$}



\end{document}
